\section{Continuous-Exponent \ScaleCert}

The lower bounds motivate a set-valued output. First consider a finite dictionary $\mathcal A=\{\alpha_1,\ldots,\alpha_L\}$ and
$\Risk_\theta(M)=\theta_0+\sum_{\ell}\theta_\ell\phi_{\alpha_\ell}(M)$ with $\theta\ge0$. Let
$z_\delta=\Phi^{-1}(1-\delta/(2m))$ and pilot bands
$[\ell_i,u_i]=[Y_i-z_\delta s_i,Y_i+z_\delta s_i]$.
Minimizing and maximizing $\Risk_\theta(M_\star)$ subject to
$\ell_i\le\Risk_\theta(M_i)\le u_i$ gives the sharp finite-dictionary interval by two linear programs.

\begin{theorem}[Finite-dictionary simultaneous coverage]
\label{thm:certificate}
If the true risk belongs to the finite dictionary and Eq.~\eqref{eq:observations} holds, then with probability at least $1-\delta$ the resulting LP interval contains $\Risk(M)$ for every target $M$ simultaneously. The interval is sharp relative to the coefficient vectors not rejected by the pilot bands.
\end{theorem}

A grid plug-in, however, does not cover off-grid exponents. We now certify the continuous class in Eq.~\eqref{eq:continuous-mixture}. Choose
$\gamma_0<\cdots<\gamma_L$ spanning $[\alpha_{\min},\alpha_{\max}]$ and let $\Delta_\ell=\gamma_{\ell+1}-\gamma_\ell$. Since
\begin{equation}
 \partial_\alpha^2\phi_\alpha(M)
 =\sum_{j>M}(\log j)^2j^{-1-\alpha}>0,
 \label{eq:basis-curvature}
\end{equation}
$\phi_\alpha(M)$ is convex in its exponent. Let $S_{\ell,\alpha}(M)$ be the secant interpolation between the two cell endpoints.

\begin{lemma}[Quadratic cell correction]
\label{lem:secant}
For $\alpha\in[\gamma_\ell,\gamma_{\ell+1}]$,
\begin{equation}
 0\le S_{\ell,\alpha}(M)-\phi_\alpha(M)
 \le \varepsilon_\ell(M):=
 \frac{\Delta_\ell^2}{8}\kappa_\ell(M),
 \label{eq:secant-bound}
\end{equation}
where $\kappa_\ell(M)$ is any upper bound on
$\sup_{\alpha\in[\gamma_\ell,\gamma_{\ell+1}]}
\partial_\alpha^2\phi_\alpha(M)$. For zeta tails we use an explicit finite-sum-plus-integral bound given in the supplement.
\end{lemma}

For the mass of $\mu$ inside cell $\ell$, define its barycentric endpoint masses $p_\ell,q_\ell\ge0$. Collect
$\vartheta=(E,p_0,q_0,\ldots,p_{L-1},q_{L-1})\ge0$. Let
\begin{align}
 D_M\vartheta&:=E+\sum_\ell\{p_\ell\phi_{\gamma_\ell}(M)
 +q_\ell\phi_{\gamma_{\ell+1}}(M)\},\\
 G_M\vartheta&:=\sum_\ell\varepsilon_\ell(M)(p_\ell+q_\ell).
 \label{eq:DG}
\end{align}
Every continuous mixture maps to such a vector and obeys
$D_M\vartheta-G_M\vartheta\le\Risk_{E,\mu}(M)\le D_M\vartheta$.
Define the outer confidence set
\begin{equation}
 \widetilde\Cset_\delta=
 \{\vartheta\ge0:D_{M_i}\vartheta\ge\ell_i,
 (D_{M_i}-G_{M_i})\vartheta\le u_i,\ i\le m\}.
 \label{eq:continuous-set}
\end{equation}
The continuous certificate is
\begin{equation}
 \widetilde L_\delta(M)=\min_{\vartheta\in\widetilde\Cset_\delta}
 (D_M-G_M)\vartheta,
 \quad
 \widetilde U_\delta(M)=\max_{\vartheta\in\widetilde\Cset_\delta}D_M\vartheta.
 \label{eq:continuous-lp}
\end{equation}
These remain linear programs.

\begin{theorem}[Off-grid simultaneous coverage]
\label{thm:continuous-certificate}
Suppose the true risk is Eq.~\eqref{eq:continuous-mixture} with a finite positive measure supported on
$[\gamma_0,\gamma_L]$, and Eq.~\eqref{eq:observations} holds. Then, with probability at least $1-\delta$, for every target $M$ simultaneously,
\begin{equation}
 \Risk_{E,\mu}(M)\in
 [\widetilde L_\delta(M),\widetilde U_\delta(M)].
 \label{eq:continuous-coverage}
\end{equation}
The conclusion holds for arbitrary off-grid exponents. If the true measure has total mass at most $W$, its secant representation differs from the true target by at most
$W\max_\ell\varepsilon_\ell(M)=O(W\Delta_{\max}^2)$.
\end{theorem}

An empty set rejects the declared class at the chosen tolerance. An unbounded upper LP is an honest abstention. Since the first coordinate of $\vartheta$ is the floor $E$, the same LPs with objectives $(D_M-G_M)\vartheta-E$ and $D_M\vartheta-E$ define excess-risk bounds $\widetilde L^{\mathrm{ex}}_\delta(M)$ and $\widetilde U^{\mathrm{ex}}_\delta(M)$. They simultaneously cover $\Risk(M)-E_{\Risk}$. Thus an excess-risk target $\epsilon$ is certified as reached when $\widetilde U^{\mathrm{ex}}_\delta(M_\star)\le\epsilon$, not reached when $\widetilde L^{\mathrm{ex}}_\delta(M_\star)>\epsilon$, and otherwise remains unidentified. Bounded mismatch is handled by expanding each pilot band by a declared radius $\rho_i$.

The same risk envelope certifies a resource requirement rather than only a loss at one chosen scale.

\begin{corollary}[Compute-to-target certificate]
\label{cor:compute-certificate}
On the simultaneous event in Theorem~\ref{thm:continuous-certificate}, define
\begin{align}
 \underline C^{\mathrm{ex}}_\delta(\epsilon)
 &:=\inf\{M\ge M_{\max}:
 \widetilde L^{\mathrm{ex}}_\delta(M)\le\epsilon\},\\
 \overline C^{\mathrm{ex}}_\delta(\epsilon)
 &:=\inf\{M\ge M_{\max}:
 \widetilde U^{\mathrm{ex}}_\delta(M)\le\epsilon\}.
 \label{eq:compute-certificate}
\end{align}
Then, for every excess-risk threshold $\epsilon>0$ simultaneously,
\begin{equation}
 \underline C^{\mathrm{ex}}_\delta(\epsilon)
 \le C^{\mathrm{ex}}_{\Risk}(\epsilon)
 \le \overline C^{\mathrm{ex}}_\delta(\epsilon).
 \label{eq:compute-coverage}
\end{equation}
Either endpoint may be infinite. Thus an unbounded upper resource requirement records that the pilot data cannot guarantee the target excess risk under the declared class.
\end{corollary}

\paragraph{Practical workflow.}
A certificate is auditable only when its structural choices are visible. We recommend: (i) declare the exponent interval, positivity assumption, and mismatch radius before inspecting the target-scale result; (ii) form simultaneous pilot bands from repeated runs or a justified error model; (iii) solve both target LPs, or invert their envelopes using Corollary~\ref{cor:compute-certificate}; and (iv) report empty sets as model rejection and unbounded sets as abstention. Comparing nested exponent intervals and mismatch radii produces a coverage--informativeness frontier rather than one opaque confidence number.

\section{Bridge to Trained Finite-Sample Regression}

The spectral curve in Eq.~\eqref{eq:spectral-tail} is a population truncation risk. It also determines the risk of a trained estimator exactly.

\begin{proposition}[Exact Gaussian OLS inflation]
\label{prop:ols}
Assume the first $M$ covariates are Gaussian and whitened, while the omitted-coordinate residual is independent and has variance $\Risk(M)$. Train ordinary least squares on $n>M+1$ independent samples, and let $\widehat\Risk_{\mathrm{OLS}}(M)$ denote the population test risk of the fitted predictor. Then
\begin{equation}
 \E\widehat\Risk_{\mathrm{OLS}}(M)
 =\Risk(M)\frac{n-1}{n-M-1}.
 \label{eq:ols-mean}
\end{equation}
Consequently,
$\frac{n-M-1}{n-1}\widehat\Risk_{\mathrm{OLS}}(M)$ is unbiased for the spectral risk at the mean level. If the omitted-coordinate residual is also Gaussian, then moreover
\begin{equation}
 \frac{\widehat\Risk_{\mathrm{OLS}}(M)}{\Risk(M)}
 \overset d=1+\frac{M}{n-M+1}F_{M,n-M+1}.
 \label{eq:ols-distribution}
\end{equation}
\end{proposition}

The omitted coordinates plus label noise form an independent residual with variance $\Risk(M)$. Equation~\eqref{eq:ols-mean} then follows from the inverse-Wishart identity. This gives an exact, inexpensive way to validate certificates on trained finite-sample learning curves rather than only on population formulas.
